\documentclass[11pt,a4paper]{article}
\newcommand{\gradelabel}{Grade 9}
\newcommand{\leconhdr}{Faire clignoter une LED}
% Préambule commun aux plans de cours de robotique (Grades 7 à 12).
% Chaque plan de cours commence par \documentclass[11pt,a4paper]{article}, définit
% \gradelabel et \leconhdr (en-têtes de page), puis % Préambule commun aux plans de cours de robotique (Grades 7 à 12).
% Chaque plan de cours commence par \documentclass[11pt,a4paper]{article}, définit
% \gradelabel et \leconhdr (en-têtes de page), puis % Préambule commun aux plans de cours de robotique (Grades 7 à 12).
% Chaque plan de cours commence par \documentclass[11pt,a4paper]{article}, définit
% \gradelabel et \leconhdr (en-têtes de page), puis \input{../../preambule_robotique.tex}.
\usepackage[T1]{fontenc}
\usepackage[utf8]{inputenc}
\usepackage{lmodern}
\usepackage[french]{babel}
\usepackage[margin=2cm,headheight=14pt]{geometry}
\usepackage[most]{tcolorbox}
\usepackage{xcolor,tabularx,booktabs,array,enumitem,fancyhdr,amssymb,tikz}
\usetikzlibrary{arrows.meta,positioning}

\pagestyle{fancy}
\fancyhf{}
\lhead{\small Robotique -- \gradelabel{} -- Novembre 2026}
\rhead{\small \leconhdr}
\cfoot{\thepage}
\setlist{itemsep=1pt,topsep=2pt,leftmargin=*}
\renewcommand{\arraystretch}{1.35}
\newcolumntype{Y}{>{\raggedright\arraybackslash}X}
\newcommand{\cb}{$\square$\,}
\newcommand{\fl}{$\rightarrow$}
\newcommand{\lig}{\noindent\rule{\linewidth}{0.3pt}\\[5mm]}

\newtcolorbox{etape}[3]{enhanced,breakable,colback=#1!6,colframe=#1!70!black,
  coltitle=white,fonttitle=\bfseries,title={#2\hfill #3},boxrule=0.8pt,arc=2mm,
  before skip=8pt,after skip=8pt}
\newtcolorbox{astuce}[1][Différenciation et besoins particuliers]{breakable,
  colback=yellow!12,colframe=yellow!50!black,fonttitle=\bfseries\small,
  title={#1},boxrule=0.5pt,arc=1mm,before skip=6pt,after skip=6pt}
\newtcolorbox{code}[1][Code]{breakable,colback=black!4,colframe=black!60,
  fonttitle=\bfseries\small,title={#1},boxrule=0.5pt,arc=1mm,fontupper=\ttfamily\small,
  before skip=6pt,after skip=6pt}
% Blocs de code coloré (écrire le code sans accents : comme dans l'éditeur)
\newtcblisting{arduino}[1][Code Arduino]{listing only,breakable,colback=black!4,
  colframe=black!60,fonttitle=\bfseries\small,title={#1},boxrule=0.5pt,arc=1mm,
  before skip=6pt,after skip=6pt,
  listing options={language=C++,basicstyle=\ttfamily\small,columns=fullflexible,
  keepspaces=true,showstringspaces=false,commentstyle=\color{green!40!black}\itshape,
  keywordstyle=\color{blue!70!black}\bfseries}}
\newtcblisting{python}[1][Code Python]{listing only,breakable,colback=black!4,
  colframe=black!60,fonttitle=\bfseries\small,title={#1},boxrule=0.5pt,arc=1mm,
  before skip=6pt,after skip=6pt,
  listing options={language=Python,basicstyle=\ttfamily\small,columns=fullflexible,
  keepspaces=true,showstringspaces=false,commentstyle=\color{green!40!black}\itshape,
  keywordstyle=\color{blue!70!black}\bfseries}}
\newcommand{\entete}[2]{%
  \begin{tcolorbox}[colback=#1!10,colframe=#1!70!black,boxrule=0.8pt,arc=2mm]
  {\Large\bfseries #2}\hfill Nom : \rule{4cm}{0.3pt}\quad Date : \rule{2.2cm}{0.3pt}
  \end{tcolorbox}}

% Bandeau de titre : #1 titre, #2 semaine
\newcommand{\titreblock}[2]{%
\begin{tcolorbox}[colback=teal!10,colframe=teal!70!black,arc=2mm,boxrule=1pt]
\centering
{\LARGE\bfseries #1}\\[2mm]
{\large Robotique -- \gradelabel{} -- Plan de cours de 2 heures}\\[1mm]
Enseignant : Francis Ki \quad|\quad #2
\end{tcolorbox}}

% Tableau du minutage : 6 contenus (accroche, mini-leçon, guidée, autonome, mise en commun, exit ticket)
\newcommand{\minutage}[6]{%
\begin{tabularx}{\linewidth}{>{\bfseries}p{2.2cm}p{2.6cm}Y p{3cm}}
\toprule
Minutes & Étape & Contenu & Organisation\\
\midrule
0 -- 10 & Accroche & #1 & Classe entière\\
10 -- 30 & Mini-leçon & #2 & Classe entière\\
30 -- 60 & Pratique guidée & #3 & Collectif puis binômes\\
60 -- 95 & Pratique autonome & #4 & Individuel sur poste\\
95 -- 110 & Mise en commun & #5 & Classe entière\\
110 -- 120 & Exit ticket & #6 & Individuel sur poste\\
\bottomrule
\end{tabularx}}

% Étayage progressif : #1 aide 1, #2 aide 2, #3 aide 3
\newcommand{\aides}[3]{%
\textbf{Étayage progressif de l'aide.}
\begin{center}
\begin{tabularx}{\linewidth}{>{\bfseries}p{2.4cm}Y p{4.4cm}}
\toprule
Niveau d'aide & Contenu & Qui intervient\\
\midrule
Aide 1 & #1 & Le binôme seul\\
Aide 2 & #2 & Le binôme avec la carte d'aide\\
Aide 3 & #3 & L'enseignant\\
\bottomrule
\end{tabularx}
\end{center}}

% Les trois groupes de la pratique autonome : #1 tâche A, #2 tâche B, #3 tâche C
\newcommand{\groupes}[3]{%
\textbf{Consigne commune\,:} chaque élève ouvre le fichier de son groupe, le renomme avec son prénom et travaille seul sur son poste. L'enseignant passe d'abord auprès du groupe A pendant 10 min.

\medskip
\begin{tabularx}{\linewidth}{>{\bfseries}p{3cm}Y p{2.2cm}}
\toprule
Groupe & Tâche sur poste & Annexe\\
\midrule
A : en dessous du niveau & #1 & Annexe C\\
B : au niveau & #2 & Annexe D\\
C : au-dessus du niveau & #3 & Annexe E\\
\bottomrule
\end{tabularx}}

% Mise en commun : #1 groupe A, #2 groupe B, #3 groupe C, #4 conclusion
\newcommand{\miseencommun}[4]{%
\begin{etape}{teal}{5. Mise en commun}{95 -- 110 min}
\begin{enumerate}
 \item (2 min) Retour au calme, écrans mis en veille.
 \item (9 min) Trois passages courts de 3 min, avec valorisation :
 \begin{itemize}
  \item \textbf{Groupe A} : #1
  \item \textbf{Groupe B} : #2
  \item \textbf{Groupe C} : #3
 \end{itemize}
 \item (4 min) Conclusion : #4
\end{enumerate}
\end{etape}}

% Clôture : #1 question 1, #2 question 2
\newcommand{\cloture}[2]{%
\begin{etape}{gray}{6. Clôture : exit ticket}{110 -- 120 min}
Formulaire en ligne ou papier (annexe G) :
\begin{enumerate}
 \item #1
 \item #2
 \item Auto-évaluation : \emph{Je maîtrise / Je progresse / J'ai besoin d'aide.}
\end{enumerate}
Lire les réponses avant la séance suivante pour ajuster les groupes de besoin.
\end{etape}}

% Grille d'observation : #1 #2 #3 intitulés de colonnes
\newcommand{\grille}[3]{%
\subsection*{Grille d'observation (enseignant)}
\begin{tabularx}{\linewidth}{Y p{2.4cm} p{2.4cm} p{2.4cm} Y}
\toprule
Élève & #1 & #2 & #3 & Groupe suivant\\
\midrule
 & & & & \\ & & & & \\ & & & & \\ & & & & \\ & & & & \\
\bottomrule
\end{tabularx}}

% Signature de l'exit ticket
\newcommand{\autoeval}{Je me sens : \cb Je maîtrise \quad \cb Je progresse \quad \cb J'ai besoin d'aide}
.
\usepackage[T1]{fontenc}
\usepackage[utf8]{inputenc}
\usepackage{lmodern}
\usepackage[french]{babel}
\usepackage[margin=2cm,headheight=14pt]{geometry}
\usepackage[most]{tcolorbox}
\usepackage{xcolor,tabularx,booktabs,array,enumitem,fancyhdr,amssymb,tikz}
\usetikzlibrary{arrows.meta,positioning}

\pagestyle{fancy}
\fancyhf{}
\lhead{\small Robotique -- \gradelabel{} -- Novembre 2026}
\rhead{\small \leconhdr}
\cfoot{\thepage}
\setlist{itemsep=1pt,topsep=2pt,leftmargin=*}
\renewcommand{\arraystretch}{1.35}
\newcolumntype{Y}{>{\raggedright\arraybackslash}X}
\newcommand{\cb}{$\square$\,}
\newcommand{\fl}{$\rightarrow$}
\newcommand{\lig}{\noindent\rule{\linewidth}{0.3pt}\\[5mm]}

\newtcolorbox{etape}[3]{enhanced,breakable,colback=#1!6,colframe=#1!70!black,
  coltitle=white,fonttitle=\bfseries,title={#2\hfill #3},boxrule=0.8pt,arc=2mm,
  before skip=8pt,after skip=8pt}
\newtcolorbox{astuce}[1][Différenciation et besoins particuliers]{breakable,
  colback=yellow!12,colframe=yellow!50!black,fonttitle=\bfseries\small,
  title={#1},boxrule=0.5pt,arc=1mm,before skip=6pt,after skip=6pt}
\newtcolorbox{code}[1][Code]{breakable,colback=black!4,colframe=black!60,
  fonttitle=\bfseries\small,title={#1},boxrule=0.5pt,arc=1mm,fontupper=\ttfamily\small,
  before skip=6pt,after skip=6pt}
% Blocs de code coloré (écrire le code sans accents : comme dans l'éditeur)
\newtcblisting{arduino}[1][Code Arduino]{listing only,breakable,colback=black!4,
  colframe=black!60,fonttitle=\bfseries\small,title={#1},boxrule=0.5pt,arc=1mm,
  before skip=6pt,after skip=6pt,
  listing options={language=C++,basicstyle=\ttfamily\small,columns=fullflexible,
  keepspaces=true,showstringspaces=false,commentstyle=\color{green!40!black}\itshape,
  keywordstyle=\color{blue!70!black}\bfseries}}
\newtcblisting{python}[1][Code Python]{listing only,breakable,colback=black!4,
  colframe=black!60,fonttitle=\bfseries\small,title={#1},boxrule=0.5pt,arc=1mm,
  before skip=6pt,after skip=6pt,
  listing options={language=Python,basicstyle=\ttfamily\small,columns=fullflexible,
  keepspaces=true,showstringspaces=false,commentstyle=\color{green!40!black}\itshape,
  keywordstyle=\color{blue!70!black}\bfseries}}
\newcommand{\entete}[2]{%
  \begin{tcolorbox}[colback=#1!10,colframe=#1!70!black,boxrule=0.8pt,arc=2mm]
  {\Large\bfseries #2}\hfill Nom : \rule{4cm}{0.3pt}\quad Date : \rule{2.2cm}{0.3pt}
  \end{tcolorbox}}

% Bandeau de titre : #1 titre, #2 semaine
\newcommand{\titreblock}[2]{%
\begin{tcolorbox}[colback=teal!10,colframe=teal!70!black,arc=2mm,boxrule=1pt]
\centering
{\LARGE\bfseries #1}\\[2mm]
{\large Robotique -- \gradelabel{} -- Plan de cours de 2 heures}\\[1mm]
Enseignant : Francis Ki \quad|\quad #2
\end{tcolorbox}}

% Tableau du minutage : 6 contenus (accroche, mini-leçon, guidée, autonome, mise en commun, exit ticket)
\newcommand{\minutage}[6]{%
\begin{tabularx}{\linewidth}{>{\bfseries}p{2.2cm}p{2.6cm}Y p{3cm}}
\toprule
Minutes & Étape & Contenu & Organisation\\
\midrule
0 -- 10 & Accroche & #1 & Classe entière\\
10 -- 30 & Mini-leçon & #2 & Classe entière\\
30 -- 60 & Pratique guidée & #3 & Collectif puis binômes\\
60 -- 95 & Pratique autonome & #4 & Individuel sur poste\\
95 -- 110 & Mise en commun & #5 & Classe entière\\
110 -- 120 & Exit ticket & #6 & Individuel sur poste\\
\bottomrule
\end{tabularx}}

% Étayage progressif : #1 aide 1, #2 aide 2, #3 aide 3
\newcommand{\aides}[3]{%
\textbf{Étayage progressif de l'aide.}
\begin{center}
\begin{tabularx}{\linewidth}{>{\bfseries}p{2.4cm}Y p{4.4cm}}
\toprule
Niveau d'aide & Contenu & Qui intervient\\
\midrule
Aide 1 & #1 & Le binôme seul\\
Aide 2 & #2 & Le binôme avec la carte d'aide\\
Aide 3 & #3 & L'enseignant\\
\bottomrule
\end{tabularx}
\end{center}}

% Les trois groupes de la pratique autonome : #1 tâche A, #2 tâche B, #3 tâche C
\newcommand{\groupes}[3]{%
\textbf{Consigne commune\,:} chaque élève ouvre le fichier de son groupe, le renomme avec son prénom et travaille seul sur son poste. L'enseignant passe d'abord auprès du groupe A pendant 10 min.

\medskip
\begin{tabularx}{\linewidth}{>{\bfseries}p{3cm}Y p{2.2cm}}
\toprule
Groupe & Tâche sur poste & Annexe\\
\midrule
A : en dessous du niveau & #1 & Annexe C\\
B : au niveau & #2 & Annexe D\\
C : au-dessus du niveau & #3 & Annexe E\\
\bottomrule
\end{tabularx}}

% Mise en commun : #1 groupe A, #2 groupe B, #3 groupe C, #4 conclusion
\newcommand{\miseencommun}[4]{%
\begin{etape}{teal}{5. Mise en commun}{95 -- 110 min}
\begin{enumerate}
 \item (2 min) Retour au calme, écrans mis en veille.
 \item (9 min) Trois passages courts de 3 min, avec valorisation :
 \begin{itemize}
  \item \textbf{Groupe A} : #1
  \item \textbf{Groupe B} : #2
  \item \textbf{Groupe C} : #3
 \end{itemize}
 \item (4 min) Conclusion : #4
\end{enumerate}
\end{etape}}

% Clôture : #1 question 1, #2 question 2
\newcommand{\cloture}[2]{%
\begin{etape}{gray}{6. Clôture : exit ticket}{110 -- 120 min}
Formulaire en ligne ou papier (annexe G) :
\begin{enumerate}
 \item #1
 \item #2
 \item Auto-évaluation : \emph{Je maîtrise / Je progresse / J'ai besoin d'aide.}
\end{enumerate}
Lire les réponses avant la séance suivante pour ajuster les groupes de besoin.
\end{etape}}

% Grille d'observation : #1 #2 #3 intitulés de colonnes
\newcommand{\grille}[3]{%
\subsection*{Grille d'observation (enseignant)}
\begin{tabularx}{\linewidth}{Y p{2.4cm} p{2.4cm} p{2.4cm} Y}
\toprule
Élève & #1 & #2 & #3 & Groupe suivant\\
\midrule
 & & & & \\ & & & & \\ & & & & \\ & & & & \\ & & & & \\
\bottomrule
\end{tabularx}}

% Signature de l'exit ticket
\newcommand{\autoeval}{Je me sens : \cb Je maîtrise \quad \cb Je progresse \quad \cb J'ai besoin d'aide}
.
\usepackage[T1]{fontenc}
\usepackage[utf8]{inputenc}
\usepackage{lmodern}
\usepackage[french]{babel}
\usepackage[margin=2cm,headheight=14pt]{geometry}
\usepackage[most]{tcolorbox}
\usepackage{xcolor,tabularx,booktabs,array,enumitem,fancyhdr,amssymb,tikz}
\usetikzlibrary{arrows.meta,positioning}

\pagestyle{fancy}
\fancyhf{}
\lhead{\small Robotique -- \gradelabel{} -- Novembre 2026}
\rhead{\small \leconhdr}
\cfoot{\thepage}
\setlist{itemsep=1pt,topsep=2pt,leftmargin=*}
\renewcommand{\arraystretch}{1.35}
\newcolumntype{Y}{>{\raggedright\arraybackslash}X}
\newcommand{\cb}{$\square$\,}
\newcommand{\fl}{$\rightarrow$}
\newcommand{\lig}{\noindent\rule{\linewidth}{0.3pt}\\[5mm]}

\newtcolorbox{etape}[3]{enhanced,breakable,colback=#1!6,colframe=#1!70!black,
  coltitle=white,fonttitle=\bfseries,title={#2\hfill #3},boxrule=0.8pt,arc=2mm,
  before skip=8pt,after skip=8pt}
\newtcolorbox{astuce}[1][Différenciation et besoins particuliers]{breakable,
  colback=yellow!12,colframe=yellow!50!black,fonttitle=\bfseries\small,
  title={#1},boxrule=0.5pt,arc=1mm,before skip=6pt,after skip=6pt}
\newtcolorbox{code}[1][Code]{breakable,colback=black!4,colframe=black!60,
  fonttitle=\bfseries\small,title={#1},boxrule=0.5pt,arc=1mm,fontupper=\ttfamily\small,
  before skip=6pt,after skip=6pt}
% Blocs de code coloré (écrire le code sans accents : comme dans l'éditeur)
\newtcblisting{arduino}[1][Code Arduino]{listing only,breakable,colback=black!4,
  colframe=black!60,fonttitle=\bfseries\small,title={#1},boxrule=0.5pt,arc=1mm,
  before skip=6pt,after skip=6pt,
  listing options={language=C++,basicstyle=\ttfamily\small,columns=fullflexible,
  keepspaces=true,showstringspaces=false,commentstyle=\color{green!40!black}\itshape,
  keywordstyle=\color{blue!70!black}\bfseries}}
\newtcblisting{python}[1][Code Python]{listing only,breakable,colback=black!4,
  colframe=black!60,fonttitle=\bfseries\small,title={#1},boxrule=0.5pt,arc=1mm,
  before skip=6pt,after skip=6pt,
  listing options={language=Python,basicstyle=\ttfamily\small,columns=fullflexible,
  keepspaces=true,showstringspaces=false,commentstyle=\color{green!40!black}\itshape,
  keywordstyle=\color{blue!70!black}\bfseries}}
\newcommand{\entete}[2]{%
  \begin{tcolorbox}[colback=#1!10,colframe=#1!70!black,boxrule=0.8pt,arc=2mm]
  {\Large\bfseries #2}\hfill Nom : \rule{4cm}{0.3pt}\quad Date : \rule{2.2cm}{0.3pt}
  \end{tcolorbox}}

% Bandeau de titre : #1 titre, #2 semaine
\newcommand{\titreblock}[2]{%
\begin{tcolorbox}[colback=teal!10,colframe=teal!70!black,arc=2mm,boxrule=1pt]
\centering
{\LARGE\bfseries #1}\\[2mm]
{\large Robotique -- \gradelabel{} -- Plan de cours de 2 heures}\\[1mm]
Enseignant : Francis Ki \quad|\quad #2
\end{tcolorbox}}

% Tableau du minutage : 6 contenus (accroche, mini-leçon, guidée, autonome, mise en commun, exit ticket)
\newcommand{\minutage}[6]{%
\begin{tabularx}{\linewidth}{>{\bfseries}p{2.2cm}p{2.6cm}Y p{3cm}}
\toprule
Minutes & Étape & Contenu & Organisation\\
\midrule
0 -- 10 & Accroche & #1 & Classe entière\\
10 -- 30 & Mini-leçon & #2 & Classe entière\\
30 -- 60 & Pratique guidée & #3 & Collectif puis binômes\\
60 -- 95 & Pratique autonome & #4 & Individuel sur poste\\
95 -- 110 & Mise en commun & #5 & Classe entière\\
110 -- 120 & Exit ticket & #6 & Individuel sur poste\\
\bottomrule
\end{tabularx}}

% Étayage progressif : #1 aide 1, #2 aide 2, #3 aide 3
\newcommand{\aides}[3]{%
\textbf{Étayage progressif de l'aide.}
\begin{center}
\begin{tabularx}{\linewidth}{>{\bfseries}p{2.4cm}Y p{4.4cm}}
\toprule
Niveau d'aide & Contenu & Qui intervient\\
\midrule
Aide 1 & #1 & Le binôme seul\\
Aide 2 & #2 & Le binôme avec la carte d'aide\\
Aide 3 & #3 & L'enseignant\\
\bottomrule
\end{tabularx}
\end{center}}

% Les trois groupes de la pratique autonome : #1 tâche A, #2 tâche B, #3 tâche C
\newcommand{\groupes}[3]{%
\textbf{Consigne commune\,:} chaque élève ouvre le fichier de son groupe, le renomme avec son prénom et travaille seul sur son poste. L'enseignant passe d'abord auprès du groupe A pendant 10 min.

\medskip
\begin{tabularx}{\linewidth}{>{\bfseries}p{3cm}Y p{2.2cm}}
\toprule
Groupe & Tâche sur poste & Annexe\\
\midrule
A : en dessous du niveau & #1 & Annexe C\\
B : au niveau & #2 & Annexe D\\
C : au-dessus du niveau & #3 & Annexe E\\
\bottomrule
\end{tabularx}}

% Mise en commun : #1 groupe A, #2 groupe B, #3 groupe C, #4 conclusion
\newcommand{\miseencommun}[4]{%
\begin{etape}{teal}{5. Mise en commun}{95 -- 110 min}
\begin{enumerate}
 \item (2 min) Retour au calme, écrans mis en veille.
 \item (9 min) Trois passages courts de 3 min, avec valorisation :
 \begin{itemize}
  \item \textbf{Groupe A} : #1
  \item \textbf{Groupe B} : #2
  \item \textbf{Groupe C} : #3
 \end{itemize}
 \item (4 min) Conclusion : #4
\end{enumerate}
\end{etape}}

% Clôture : #1 question 1, #2 question 2
\newcommand{\cloture}[2]{%
\begin{etape}{gray}{6. Clôture : exit ticket}{110 -- 120 min}
Formulaire en ligne ou papier (annexe G) :
\begin{enumerate}
 \item #1
 \item #2
 \item Auto-évaluation : \emph{Je maîtrise / Je progresse / J'ai besoin d'aide.}
\end{enumerate}
Lire les réponses avant la séance suivante pour ajuster les groupes de besoin.
\end{etape}}

% Grille d'observation : #1 #2 #3 intitulés de colonnes
\newcommand{\grille}[3]{%
\subsection*{Grille d'observation (enseignant)}
\begin{tabularx}{\linewidth}{Y p{2.4cm} p{2.4cm} p{2.4cm} Y}
\toprule
Élève & #1 & #2 & #3 & Groupe suivant\\
\midrule
 & & & & \\ & & & & \\ & & & & \\ & & & & \\ & & & & \\
\bottomrule
\end{tabularx}}

% Signature de l'exit ticket
\newcommand{\autoeval}{Je me sens : \cb Je maîtrise \quad \cb Je progresse \quad \cb J'ai besoin d'aide}


\begin{document}
\titreblock{Faire clignoter une LED}{Semaine du 2 au 6 novembre 2026}

%==================================================================
\section{Informations générales}
\begin{tabularx}{\linewidth}{>{\bfseries}p{3.6cm}Y}
\toprule
Titre & Faire clignoter une LED avec Arduino\\
Classe / effectif & Grade 9, 11 à 20 élèves, niveaux hétérogènes, élèves à besoins particuliers\\
Durée & 2 heures (120 min)\\
Lieu & Salle informatique, un poste par élève\\
Matériel et logiciels &
\begin{itemize}
 \item Un poste par élève avec navigateur web : \textbf{Tinkercad Circuits} (tinkercad.com, gratuit, compte élève ou code de classe) pour construire et simuler le montage.
 \item Arduino IDE installé sur les postes (vu en septembre) pour copier, vérifier et éventuellement téléverser le code.
 \item Facultatif : une carte Arduino Uno avec LED et résistance de 220~$\Omega$ par binôme pour tester le montage réel.
 \item Vidéoprojecteur ou tableau interactif.
\end{itemize}\\
Pré-requis &
Séances précédentes : carte Arduino, entrées et sorties, structure d'un programme (\texttt{setup} et \texttt{loop}), premier montage d'une LED (octobre).\\
\bottomrule
\end{tabularx}

\begin{astuce}[Préparation avant le cours (environ 40 min)]
\begin{enumerate}
 \item Créer une classe dans Tinkercad (si disponible) ou un code d'invitation, et vérifier que chaque élève se connecte.
 \item Préparer un circuit modèle partagé (Arduino Uno, LED, résistance 220~$\Omega$ sur la broche 8) que les élèves dupliquent.
 \item Préparer les fichiers des annexes C à E et le formulaire de sortie (annexe G).
 \item Si des cartes réelles sont utilisées, vérifier les câbles USB et le port dans l'IDE.
\end{enumerate}
\end{astuce}

\begin{astuce}[Sécurité et bonnes pratiques]
Une LED se branche toujours avec une résistance. Pas de court-circuit entre la broche et la masse. Débrancher la carte avant de modifier le montage réel.
\end{astuce}

%==================================================================
\section{Objectifs d'apprentissage}
\begin{tcolorbox}[colback=gray!8,colframe=gray!60!black,arc=1mm,boxrule=0.5pt]
\textbf{Objectifs du programme}
\begin{itemize}
 \item Utiliser les temporisations.
 \item Modifier un programme existant.
\end{itemize}
\textbf{À la fin de la séance, l'élève est capable\,:}
\begin{itemize}
 \item d'expliquer le rôle de \texttt{pinMode}, \texttt{digitalWrite} et \texttt{delay}\,;
 \item de faire clignoter une LED à une durée donnée et de calculer la période et la fréquence\,;
 \item de modifier un programme existant pour changer le rythme, et de vérifier le résultat en simulation.
\end{itemize}
\textbf{Compétences numériques\,:} construire un circuit en simulation, démarrer et arrêter la simulation, copier du code vers l'IDE.
\end{tcolorbox}

%==================================================================
\section{Déroulé minuté (120 min)}
\minutage{Le feu clignotant et le battement de c\oe ur}{setup, loop, delay et fréquence}{Un clignotement à 1 s en collectif, puis un rythme modifié en binômes}{Trois niveaux de besoin}{Valoriser chaque groupe}{Trois questions sur poste}

%------------------------------------------------------------------
\begin{etape}{orange}{1. Accroche : le clignotant et le battement de c\oe ur}{0 -- 10 min}
\begin{enumerate}
 \item (5 min) Projeter une vidéo courte ou un schéma d'un clignotant de voiture, d'un phare et d'un moniteur cardiaque. \emph{Comment sont-ils commandés ? Combien de fois par seconde ?}
 \item (5 min) Les élèves tapent avec la main au rythme d'un clignotant (1 fois par seconde), puis plus vite. Question-problème : \textbf{«~Comment dire à une carte électronique de s'allumer et de s'éteindre à un rythme choisi ?~»}
\end{enumerate}
\end{etape}

%------------------------------------------------------------------
\begin{etape}{blue}{2. Mini-leçon : les instructions de base}{10 -- 30 min}
\textbf{Support\,:} annexe A. Quatre temps de 5 min.

\textbf{Temps 1 : le montage.} La LED est branchée en série avec une résistance sur une broche numérique (broche 8) et à la masse (GND).

\textbf{Temps 2 : le programme.}
\begin{arduino}[Le programme Blink]
const int LED = 8;      // numero de la broche

void setup() {
  pinMode(LED, OUTPUT); // la broche est une sortie
}

void loop() {
  digitalWrite(LED, HIGH); // allume
  delay(1000);             // attend 1000 ms
  digitalWrite(LED, LOW);  // eteint
  delay(1000);
}
\end{arduino}

\textbf{Temps 3 : le vocabulaire.}
\begin{center}
\begin{tabularx}{\linewidth}{>{\ttfamily}p{3.6cm}Y}
\toprule
Instruction & Rôle\\
\midrule
pinMode(LED, OUTPUT) & Déclare la broche comme sortie (dans \texttt{setup})\\
digitalWrite(LED, HIGH) & Met la broche à 5~V : la LED s'allume\\
digitalWrite(LED, LOW) & Met la broche à 0~V : la LED s'éteint\\
delay(1000) & Attend 1000 ms, soit 1 seconde\\
\bottomrule
\end{tabularx}
\end{center}

\textbf{Temps 4 : période et fréquence.} La \emph{période} est la durée d'un cycle (allumée + éteinte). La \emph{fréquence} est le nombre de cycles par seconde : $f = 1 \div T$. Ici $T = 2$~s, donc $f = 0{,}5$~Hz.

\begin{astuce}
Colorer les instructions de la même façon que dans l'IDE. Lire le programme à voix haute ligne par ligne. Donner un modèle avec les accolades déjà placées. Rappeler que chaque ligne se termine par un point-virgule.
\end{astuce}
\end{etape}

%------------------------------------------------------------------
\begin{etape}{green}{3. Pratique guidée}{30 -- 60 min}
\textbf{Fiche\,:} annexe B.

\textbf{3a. Clignotement à 1 s en collectif (12 min).} Dupliquer le circuit modèle, ouvrir l'éditeur de code (mode texte), saisir le programme Blink, démarrer la simulation. Chaque élève chronomètre trois clignotements.

\textbf{3b. Un rythme modifié en binômes (18 min).} Modifier uniquement les valeurs de \texttt{delay} pour obtenir un clignotement à 2 Hz, puis un battement de c\oe ur (deux éclairs rapides, une pause). Vérifier en simulation. Alternance des rôles (pilote, vérificateur).

\aides{Le programme du 3a et le tableau période et fréquence}{Questions guides : \emph{Quelle valeur règle le temps allumé ? Quelle valeur règle le temps éteint ?}}{L'enseignant modifie une valeur avec le binôme, qui poursuit}

\begin{astuce}
En cas d'erreur, lire le message de l'éditeur à voix haute et chercher d'abord un point-virgule manquant ou une accolade oubliée. Garder un exemple correct visible.
\end{astuce}
\end{etape}

%------------------------------------------------------------------
\begin{etape}{violet}{4. Pratique autonome différenciée}{60 -- 95 min}
\groupes{Programme Blink à compléter (valeurs de \texttt{delay} manquantes) avec une banque de valeurs : 200, 500, 1000, 2000. Obtenir un clignotement rapide, puis lent, et noter le résultat.}{Obtenir un clignotement à 2 Hz puis le signal SOS en morse (trois éclairs courts, trois longs, trois courts) en modifiant le programme.}{Écrire une fonction \texttt{void clignote(int duree)} qui allume et éteint la LED, puis l'utiliser pour produire un battement de c\oe ur, et tester avec la variable \texttt{duree} modifiée. Commenter chaque ligne du code.}

\medskip
\textbf{Plan B si la simulation ne fonctionne pas\,:} utiliser le programme dans l'IDE avec la vérification de compilation, et relire à voix haute avec un camarade.
\end{etape}

%------------------------------------------------------------------
\miseencommun{un élève montre son clignotement lent et rapide et explique la valeur modifiée. Valoriser le lien entre la valeur et le rythme.}{un binôme fait tourner le SOS en simulation. Valoriser l'exactitude du rythme.}{un élève présente sa fonction \texttt{clignote} et explique à quoi sert le paramètre. Valoriser la structuration.}{\textit{delay règle le rythme, la période est la somme du temps allumé et du temps éteint.}}

\cloture{Quel est le rôle de l'instruction \texttt{delay(500)} ?}{Une LED est allumée 300 ms et éteinte 700 ms : quelle est la période, quelle est la fréquence ?}

%==================================================================
% ANNEXES
%==================================================================
\newpage
\section*{Annexe A : Support résumé de la mini-leçon}
\entete{blue}{Faire clignoter une LED}

\begin{tcolorbox}[colback=yellow!10,colframe=yellow!50!black]
\textbf{Montage :} broche 8 $\rightarrow$ résistance 220~$\Omega$ $\rightarrow$ LED $\rightarrow$ GND.\\
\textbf{Période} $T$ = temps allumé + temps éteint. \textbf{Fréquence} $f = 1 \div T$ (en Hz).
\end{tcolorbox}

\begin{arduino}[Modèle]
void setup() {
  pinMode(8, OUTPUT);
}
void loop() {
  digitalWrite(8, HIGH);
  delay(1000);
  digitalWrite(8, LOW);
  delay(1000);
}
\end{arduino}

\begin{tabularx}{\linewidth}{>{\ttfamily}p{3.6cm}Y}
\toprule
Instruction & Rôle\\
\midrule
pinMode & Déclare une broche\\
digitalWrite & Écrit HIGH ou LOW\\
delay & Attend en millisecondes (1000 ms = 1 s)\\
\bottomrule
\end{tabularx}

%==================================================================
\newpage
\section*{Annexe B : Fiche de pratique guidée}
\entete{green}{Pratique guidée : clignotement}

\subsection*{3a. Mesure en collectif}
Durée de trois clignotements : \dotfill{} s \quad Période : \dotfill{} s \quad Fréquence : \dotfill{} Hz

\subsection*{3b. En binôme : changer le rythme}
\begin{tabularx}{\linewidth}{>{\bfseries}p{3.8cm}Y Y Y}
\toprule
Rythme demandé & delay allumé & delay éteint & Vérifié en simulation ?\\
\midrule
1 Hz (période 1 s) & & & \\
2 Hz (période 0,5 s) & & & \\
Battement de c\oe ur & & & \\
\bottomrule
\end{tabularx}

\begin{tcolorbox}[colback=yellow!10,colframe=yellow!50!black,title={Cartes d'aide à utiliser dans l'ordre},fonttitle=\bfseries\small]
\textbf{Aide 1 :} regarde le programme du 3a et la formule de la fréquence.\\
\textbf{Aide 2 :} Quelle valeur règle le temps allumé ? Quelle valeur règle le temps éteint ?\\
\textbf{Aide 3 :} demande à l'enseignant de modifier une valeur avec toi.
\end{tcolorbox}

%==================================================================
\newpage
\section*{Annexe C : Niveau A (en dessous du niveau)}
\entete{orange}{Niveau A : je complète le programme}

\textbf{Banque de valeurs :} 200, 500, 1000, 2000.

\begin{arduino}[Programme à compléter]
void setup() {
  pinMode(8, OUTPUT);
}
void loop() {
  digitalWrite(8, HIGH);
  delay(_____);
  digitalWrite(8, LOW);
  delay(_____);
}
\end{arduino}

\subsection*{Activité 1 : essaie les deux rythmes}
\begin{tabularx}{\linewidth}{>{\bfseries}p{3.2cm}Y Y Y}
\toprule
Rythme & Valeur delay & La LED clignote\ldots & Vrai ?\\
\midrule
Rapide & & \cb vite \quad \cb lentement & \\
Lent & & \cb vite \quad \cb lentement & \\
\bottomrule
\end{tabularx}

\subsection*{Activité 2 : vrai ou faux ?}
1. \texttt{delay(1000)} attend 1 seconde. \cb Vrai \cb Faux\\
2. \texttt{LOW} allume la LED. \cb Vrai \cb Faux

%==================================================================
\newpage
\section*{Annexe D : Niveau B (au niveau)}
\entete{blue}{Niveau B : 2 Hz et SOS}

\subsection*{Mission 1 : clignotement à 2 Hz}
Période : \dotfill{} ms \quad delay allumé : \dotfill{} \quad delay éteint : \dotfill

\subsection*{Mission 2 : SOS en morse}
Trois éclairs courts (200 ms), trois longs (600 ms), trois courts, avec 200 ms de pause entre les éclairs et une pause de 1 s à la fin.

\begin{tabularx}{\linewidth}{>{\bfseries}p{3cm}Y Y}
\toprule
Partie & Durée allumée (ms) & Durée éteinte (ms)\\
\midrule
S (court) $\times$ 3 & & \\
O (long) $\times$ 3 & & \\
S (court) $\times$ 3 & & \\
\bottomrule
\end{tabularx}

Combien d'instructions \texttt{delay} dans ton programme ? \dotfill

%==================================================================
\newpage
\section*{Annexe E : Niveau C (au-dessus du niveau)}
\entete{violet}{Niveau C : une fonction qui clignote}

\textbf{Mission :} écrire une fonction pour éviter de répéter les mêmes lignes.

\begin{arduino}[Modèle de départ]
void clignote(int duree) {
  digitalWrite(8, HIGH);
  delay(duree);
  digitalWrite(8, LOW);
  delay(duree);
}
\end{arduino}

\begin{enumerate}
 \item Utilise \texttt{clignote(200)} et \texttt{clignote(600)} dans \texttt{loop} pour créer le SOS en moins de 10 lignes.\\ \lig
 \item Crée un battement de c\oe ur (deux éclairs de 100 ms, pause de 800 ms) et explique l'intérêt d'une fonction.\\ \lig \lig
\end{enumerate}

\textbf{Défi bonus :} une variable \texttt{vitesse} qui double la fréquence quand on la change.

%==================================================================
\newpage
\section*{Annexe F : Corrections et attendus (enseignant)}

\subsection*{Annexe B}
1 Hz : 500 ms allumé et 500 ms éteint (ou 1000 et 1000 donne 0,5 Hz). 2 Hz : période 500 ms, soit 250 ms allumé et 250 ms éteint. Battement : deux éclairs (par exemple 100 ms allumé, 100 ms éteint, répété deux fois) puis 800 ms d'attente.

\subsection*{Annexe C}
Valeurs rapides : 200 ; lentes : 2000 ; la LED clignote vite pour 200, lentement pour 2000. Vrai, Faux.

\subsection*{Annexe D}
2 Hz : période 500 ms, 250 ms et 250 ms. SOS : court 200/200 ($\times$ 3), long 600/200 ($\times$ 3), court 200/200 ($\times$ 3), pause 1000 ms. Total d'au moins 18 \texttt{delay}.

\subsection*{Annexe E}
\texttt{loop} : appeler \texttt{clignote(200)} trois fois, \texttt{clignote(600)} trois fois, \texttt{clignote(200)} trois fois, \texttt{delay(1000)}. Une fonction évite de répéter et simplifie les corrections.\\
\textbf{Critères de réussite :} fonction correcte, appel avec paramètre, code commenté.

%==================================================================
\newpage
\section*{Annexe G : Exit ticket et grille d'observation}
\entete{gray}{Exit ticket : clignoter une LED}

\begin{enumerate}
 \item Quel est le rôle de l'instruction \texttt{delay(500)} ?\\ \lig
 \item Une LED est allumée 300 ms et éteinte 700 ms : quelle est la période, quelle est la fréquence ?\\ \lig
 \item \autoeval
\end{enumerate}
\textbf{Attendus :} 1. Attendre 500 ms (0,5 s). 2. Période 1000 ms $=$ 1 s, fréquence 1 Hz.

\bigskip
\grille{Montage correct}{Programme fonctionnel}{Période calculée}

\end{document}
